% PIBA Template -- Software Architecture Lab, WS 2026/27, University of Stuttgart
%
% Self-contained: compile with pdflatex (no additional files needed).
% Rules: the whole PIBA must fit on ONE page. Do not change font size or margins.
% Replace all text in \todo{...} and remove the hints before submitting.

\documentclass[11pt,a4paper]{article}

\usepackage[utf8]{inputenc}
\usepackage[T1]{fontenc}
\usepackage{lmodern}
\usepackage[a4paper,margin=2cm]{geometry}
\usepackage{enumitem}
\usepackage{graphicx}
\usepackage{xcolor}
\usepackage[hidelinks]{hyperref}

\pagestyle{empty}
\setlength{\parindent}{0pt}
\setlength{\parskip}{0.5em}
\setlist{itemsep=0.1em, topsep=0.2em}

% Placeholder marker; leave the definition, but remove all uses before submitting.
\newcommand{\todo}[1]{\textcolor{gray}{[#1]}}

% Part heading
\newcommand{\pibapart}[1]{\vspace{0.4em}{\large\textbf{#1}}\par}

% ---------------------------------------------------------------------------
% Project information
% ---------------------------------------------------------------------------
\newcommand{\projecttitle}{MealSwipe: Health-Aware Meal Planning}
\newcommand{\teamname}{Team Berlin}
\newcommand{\teammembers}{Jan, Paul, Pascal, Jenny, Georg, George}

\begin{document}

\begin{center}
  {\LARGE\textbf{\projecttitle}}\\[0.4em]
  {\large PIBA -- Software Architecture Lab, WS 2026/27}\\[0.3em]
  \teamname\ \textbullet\ \teammembers
\end{center}

% ---------------------------------------------------------------------------
\pibapart{Problem}
% What is wrong, missing, or hard today, and for whom?
Deciding what to cook every week is tedious, so people repeat the same few dishes or buy on impulse.
Recipe platforms only help if you already know what you want and ignore budget, time, or allergies.
Health ratings like the Nutri-Score are one-size-fits-all: a 900\,kcal high-protein meal suits someone building muscle but not someone trying to lose weight.

% ---------------------------------------------------------------------------
\pibapart{Idea}
% What is your proposed solution, and what does its architecture look like?
% The architecture is the heart of your idea.
A web app that finds meals by swiping: a few coarse questions (``Noodles today?''), then dish cards swiped left or right.
From the liked dishes it builds a plan of $X$ meals, weighing price, cooking time, and fit with personal goals (e.g.\ high protein), plus a priced shopping list.
\begin{itemize}
  \item \textbf{Web Client} (Flask/Jinja, plain JS for swiping): swipe UI, profile, plan, list.
  \item \textbf{Preference Engine}: holds the swipe session and picks the next question or card.
  \item \textbf{Recipe Catalog}: recipes, nutrition, and prices, filled via import adapters (mock data first).
  \item \textbf{User Profile}: goals, hard filters (vegan, allergies), and history for variety.
  \item \textbf{Health Evaluator}: rates dishes \emph{relative to the user's goals}; a rule-based and an AI (LLM) implementation with personalized explanations share one interface.
  \item \textbf{Meal Planner}: ranks candidates, selects $X$ meals, aggregates the shopping list.
\end{itemize}
\textbf{Style:} modular monolith with ports and adapters (replaceable data sources and AI provider).
\textbf{Quality attributes:} \emph{modifiability} (swap evaluators or sources), \emph{availability} (rule-based fallback if the AI fails), \emph{privacy} (only anonymized nutrient data reaches the LLM), \emph{responsiveness} (next card $<$\,200\,ms).
\textbf{Out of scope:} live prices, ordering, medical diets, native apps.

% ---------------------------------------------------------------------------
\pibapart{Benefits}
% What improves if the idea works, and how would you show it?
% Link each benefit back to the problem.
\begin{itemize}
  \item \textbf{Reduced planning effort} -- addresses tedious meal selection; shown by the time to create a weekly meal plan and shopping list vs.\ manual planning (user study).
  \item \textbf{Budget control} -- addresses impulse buys; costs are known before shopping, shown by the share of generated plans within the user's budget.
  \item \textbf{Healthier, goal-fitting diet} -- addresses one-size-fits-all ratings; shown by one dish rated differently per goal profile and plans meeting calorie/protein targets within 10\,\%.
  \item \textbf{More variety and satisfaction} -- addresses repetitive meals; shown by distinct dishes per user across sessions and participants' ratings of their generated plans.
\end{itemize}

% ---------------------------------------------------------------------------
\pibapart{Action}
% What will you build and evaluate, in which steps, within 3 months?
% Mark the minimum deliverable and keep some buffer.
\begin{enumerate}
  \item \emph{Wk 1--3, Design:} user stories, quality scenarios, interfaces, ADRs; mock data (100 recipes).
  \item \emph{Wk 4--7, Core:} catalog, profile, swipe UI, planner, rule-based evaluator.\\
  \textbf{Minimum deliverable:} swipe to $X$ meals with priced shopping list and health rating.
  \item \emph{Wk 8--10, Extend:} LLM evaluator, coarse questions, recommendations, gamification.
  \item \emph{Wk 11--12, Evaluate:} user study, quality scenario checks; buffer.
\end{enumerate}
\textbf{Tech:} Python/Flask, PostgreSQL, Docker, LLM API.
\textbf{Team:} 2 frontend, 2 backend, 2 health/AI.

% ---------------------------------------------------------------------------
\vfill
{\small\textbf{AI usage statement.} Claude (Anthropic) was used to restructure and rephrase our draft.}

\end{document}
